\section{The actual full character matrix under an explicit half-plane hypothesis}
\label{sec:character-common-matrix}

For a nontrivial character the original common-sum sector is
diagonal and has no polar modes. We construct its entire
smoothed matrix on the inherited narrow Hilbert completion,
using the actual local source and both dual inverses. The
construction has an explicit zero-free half-plane premise.
Its positive-ideal restriction is the operator already
evaluated in MA-MO-017, while its quotient action retains
the negative convolution contribution.

Let \(K\) be a fixed number field, and let \(\eta\) be a
nontrivial primitive finite-order Hecke character. Choose the
actual coherent additive character, self-dual measures, and
norm splitting of MA-NR/MA-CC/MA-CQ. The finite place set
\(S\) contains the archimedean places, all ramification of
\(\eta\), and all places required by the fixed additive
character. In addition it is \emph{structural} in Meyer's
precise sense:
\[
      C_S\simeq C_K/\mathcal O_{\overline S}^{\times}.
 \tag{MA-MO.470}
\]
Here \(\mathcal O_{\overline S}^{\times}\) is the product
of the local unit groups outside \(S\). Thus the chosen
global compact character sector is the actual corresponding
sector of \(\mathcal T(C_S)\), not an arbitrary radial
representative. Such sets exist by finiteness of the ideal
class group. This requirement is stronger than merely
containing the conductor primes.

Use probability Haar on the compact norm-one quotient and
the same centered conventions as MA-MO-012. Put
\[
 \begin{gathered}
 H=H_{\Sigma,a,s},\quad 0\le a_*<a<1/2,\quad s\ge0,\\
 I_\eta=I_\eta^c=\vartheta_\eta\mathcal T_\eta\subset\mathcal E,
 \qquad J_\eta=\overline{I_\eta}^{\,H},\qquad Q_\eta=H/J_\eta.
 \end{gathered}
 \tag{MA-MO.471}
\]
The same norm spaces in the dual compact sector are understood
when \(\bar\eta\) appears. Assume the explicit premise
\[
 L_K(\sigma+it,\omega)\ne0
       \quad\text{for }\sigma>1/2+a_*,
             \qquad \omega\in\{\eta,\bar\eta\}.
 \tag{MA-MO.472}
\]
The finite functions are entire for these nontrivial characters.
The premise concerns their positive-half-plane zeros, which
there coincide with the completed nontrivial zeros.
Let \(\phi\) be the real smooth source cutoff, equal to one
sufficiently far left and zero sufficiently far right, with
\(\phi'\in C_c^\infty\); no symmetry is imposed.

\paragraph{MA-MO-022 (full common-space character operator and its trace).}
The original smoothed common-space matrix
\(C_g(P_+-P_-)\), on its actual nontrivial smooth sector,
extends uniquely to a trace-class endomorphism
\[
          \mathscr D_{\eta,g}:H\longrightarrow H
                    \qquad(g\in C_c^\infty(\mathbb R)).
 \tag{MA-MO.473}
\]
It preserves \(J_\eta\), its restriction there is the
actual \(B_{\eta,g}\) of MA-MO-012/017, and its induced
operator on \(Q_\eta\) is exactly \(-C_g\). In particular
\[
 \begin{aligned}
 \operatorname{Tr}_H\mathscr D_{\eta,g}
   &=\tau_{\eta,S}(g,\phi)
                         -\operatorname{Tr}_{Q_\eta}C_g,\\
 \det_H(1-z\mathscr D_{\eta,g})
   &=\det_{J_\eta}(1-zB_{\eta,g})
                         \det_{Q_\eta}(1+zC_g),
                   \qquad z\in\mathbb C.
 \end{aligned}
 \tag{MA-MO.474}
\]
The local trace \(\tau_{\eta,S}\) is exactly (MA-MO.361),
including the actual finite Fourier, conductor and different
factors in \(m_\eta\), and the nonsymmetric-cutoff residual.

For the requested \(g(0)=0\) tests the full matrix trace is
\[
 \boxed{\begin{aligned}
 \operatorname{Tr}_H\mathscr D_{\eta,g}
  &=\sum_{\substack{\mathfrak p\notin S\\k\ge1}}
       \frac{\log N\mathfrak p}{(N\mathfrak p)^{k/2}}
       \bigl[c_{\mathfrak p}^{\,k}g(k\log N\mathfrak p)\\
  &\hspace{10em}
       +\overline{c_{\mathfrak p}}^{\,k}g(-k\log N\mathfrak p)\bigr].
 \end{aligned}}
 \tag{MA-MO.475}
\]
where \(c_{\mathfrak p}=\eta(\mathfrak p)\) at unramified
places. The sum is finite. No assumption that \(S\) contains
all primes contributing to the test has been made.

\begin{proof}
\emph{The actual structural sector and its smooth common space.}
Condition (MA-MO.470) is Meyer's sufficiently-large condition,
supplied \texttt{main.tex} lines~3592--3602. His actual
global/local coinvariant theorem, lines~3671--3707, identifies
the specified sector with \(\mathcal T_\eta\); it applies
without a class-number-one or unit-rank restriction. The
norm-one quotient is compact, so the norm map is proper.
On a fixed compact character, weighted Schwartz functions
are therefore measured by the one real norm coordinate.
The actual inclusion
\[
        \mathcal E_\omega
          =\mathscr S(C_S)_{\cup,\omega}
                          \longrightarrow\mathcal T_\omega,
                   \qquad\omega=\eta,\bar\eta,
 \tag{MA-MO.476}
\]
is the inclusion following the split sequence (24),
\texttt{eq:Twist\_convariants}. It is an inclusion into
the local coinvariant source topology, not a choice of
additive representative.

The singular two-dimensional Poisson representation has
trivial compact type. A nontrivial finite-order global
character is trivial on the connected norm splitting and
nontrivial on the compact quotient. By (MA-MO.470) the
specified local sector is that same nontrivial compact type.
Consequently its singular Poisson term is zero. In this
sector the original common-sum space is exactly
\(i_-\mathcal E=\{(y,y):y\in\mathcal E\}\), and
\(P_-=1\) on it. Its inherited completion is the diagonal
sum-norm space \(H\) in (MA-MO.471).

\emph{Original wide inverses on every smooth diagonal input.}
For \(\omega=\eta,\bar\eta\), write
\[
 \begin{aligned}
 A_\omega(z)&=E_{\omega,S}(1/2+z)L_K(1/2+z,\omega),\\
 E_{\omega,S}(s)&=\prod_{\substack{\mathfrak p\in S\ \mathrm{finite}\\
                                  \omega\ \mathrm{unramified}\ }}
                         (1-c_{\omega,\mathfrak p}(N\mathfrak p)^{-s}),\\
 F_{\omega,y}(x)&=
   \sum_{\mathfrak a\ \mathrm{prime\ to}\ S}
        \mu_K(\mathfrak a)c_{\omega,\mathfrak a}
        (N\mathfrak a)^{-1/2}y(x+\log N\mathfrak a).
 \end{aligned}
 \tag{MA-MO.477}
\]
These are the exact original Euler inverse coefficients of
MA-CC.1. At a wide weight \(b>1/2\) the series converges
absolutely in operator norm, dominated by
\(\zeta_K(b+1/2)\). On \(y\in\mathcal E\) the series
and every derivative converge locally absolutely, using
arbitrarily rapid right decay. Its Mellin transform on
the wide line is \(\mathcal My(z)/A_\omega(z)\).

All zeros of the finite Euler factors in (MA-MO.477) have
\(\Re z=-1/2\), so they do not affect the positive
inverse lines. Hypothesis (MA-MO.472) gives a fixed margin
to every global zero on every fixed strip \(\Re z\ge a\).
The reciprocal proof of MA-RS-004, in its fixed-Hecke version
used in MA-MO-012, gives polynomial bounds for
\(A_\omega^{-1}\) and all derivatives there. There is
no pole of a nontrivial finite \(L\)-function at one.
The entire numerator \(\mathcal My\) is rapidly decreasing
on each fixed strip. Shifting from the original wide line
to any \(b\ge a\) therefore proves
\[
 \begin{gathered}
 e^{bx}F_{\omega,y}\in\mathscr S\quad(b\ge a),\qquad
 F_{\omega,y}=J_{\omega,a}y,\\
 V_aF_{\omega,y}=j_{\omega,a}V_ay,\qquad
 j_{\omega,a}(u)=A_\omega(a+iu)^{-1}.
 \end{gathered}
 \tag{MA-MO.478}
\]
This holds for \emph{every} \(y\in\mathcal E\), not
merely ideal vectors. Equality with the original series is
proved on the wide line before the shift. No inverse is
moved through zeros in the interior narrow strip.

Put \(q=1-\phi\). The function \(qF_{\omega,y}\) belongs
to \(\mathcal E_\omega\): it is zero sufficiently far
left and has all the right weighted derivative estimates
in (MA-MO.478). By (MA-MO.476) it is consequently an actual
local source vector. This is the step that justifies the
next Fourier transform on its original domain.

\emph{Both Fourier sectors and the original positive output.}
Let \(\mathfrak S_\eta:\mathcal T_\eta\to
\mathcal T_{\bar\eta}\) and
\(\mathfrak S_{\bar\eta}:\mathcal T_{\bar\eta}\to
\mathcal T_\eta\) be the actual mutually inverse source
Fourier maps of MA-MO-017/019. They include the finite
phases, conductor and different; in particular no false
archimedean odd-parity involution is used. The original
map \(p_+\) on the smooth diagonal is
\[
 \begin{aligned}
 H_{\eta,y}&=qF_{\eta,y}
                 +\mathfrak S_{\bar\eta}(qF_{\bar\eta,Ry})
                                      \in\mathcal T_\eta,\\
 \mathfrak S_\eta H_{\eta,y}&=H_{\bar\eta,Ry},\\
 (P_+-P_-)i_-y
    &=i_-\bigl(\vartheta_\eta H_{\eta,y}-y\bigr).
 \end{aligned}
 \tag{MA-MO.479}
\]
The first line is the defining source formula with its two
wide inverses; (MA-MO.478) has identified their values.
The second line follows by applying the actual inverse
Fourier pair. Nontrivial-sector Poisson summation makes
the two global output components equal, proving the last
line. In particular the positive output is in the actual
ideal \(I_\eta\).

For any \(F\in\mathcal T_\omega\), the first source
family has weighted Schwartz estimates for every
\(c>-1/2\). Its Mellin transform is holomorphic there.
The direct theta transform is \(A_\omega\mathcal MF\)
on a wide line; shifting it to \(a\) crosses no pole,
because \(A_\omega\) is entire. Thus
\[
                  \vartheta_\omega F=Z_{\omega,a}F
                     \quad\text{on the actual local source}.
 \tag{MA-MO.480}
\]
The source output lies in \(\mathcal E\) by its two
Poisson-related tails. This direct summation step uses
no zero-free assertion: zeros are not poles of the forward
multiplier. There is no rational principal polar subtraction
or hidden moment condition in (MA-MO.480).

\emph{The positive-line crossed kernel.}
Use the actual equivariant multiplier
\(\mathfrak U_\eta=\mathfrak S_{\bar\eta}R\) and
\(m_\eta\) from MA-MO.281. Thus
\(\mathcal M\mathfrak S_{\bar\eta}F(z)
       =m_\eta(z)\mathcal MF(-z)\).
Set
\[
       z=a+it,\qquad w=a+iu,\qquad
       G(z)=\mathcal Mg(z),\qquad
       P(\xi)=\int e^{\xi x}\phi'(x)\,dx.
\]
The operator
\(C_g Z_{\eta,a}\mathfrak S_{\bar\eta}M_qJ_{\bar\eta,a}\)
has, on its original smooth core, the kernel
\[
       o_{\eta,a,g}(t,u)
          =\frac{G(z)A_\eta(z)m_\eta(z)j_{\bar\eta,a}(u)}{2\pi}
                           \frac{P(-z-w)}{-z-w}
 \tag{MA-MO.481}
\]
between the positive-line dual input and positive-line
original output. Indeed Fourier inversion on the input
line and \(\int q(x)e^{\xi x}dx=P(\xi)/\xi\) for
\(\Re\xi<0\) give exactly this formula. Absolute
justification on the core follows from its Schwartz
input transform and \(\Re(-z-w)=-2a\). The source
cutoff vector and its actual Fourier transform have the
required domains by (MA-MO.476)--(MA-MO.480).

All derivatives of (MA-MO.481) are Schwartz: the local
factor \(m_\eta\), the forward multiplier and the
dual inverse have polynomial derivative bounds on these
fixed lines; \(G(a+it)\) and \(P(-2a-i(t+u))\) give
the two rapid directions, and the denominator is separated
from zero. The Bessel/Hermite proof therefore gives an
actual extension
\(O_{\eta,a}(g):H_{a,s}^{\bar\eta}\to H_{a,s}^\eta\)
in trace class, for every \(s\ge0\). Interchanging the
characters gives the other crossed kernel. The input
factor in (MA-MO.481) is the \emph{dual} inverse; replacing
it by the same-character inverse is not permissible.

\emph{Compatible rows in the exact sum norm.}
Let \(K_{\eta,a}(g)\) be the one-sided commutator kernel
(MA-MO.184). Under (MA-MO.472) it agrees on all of
\(\mathcal E\) with
\(C_g(M_\phi-Z_{\eta,a}M_\phi J_{\eta,a})\), by
(MA-MO.478) and the graph cutoff identity
\(j_{\eta,a}\chi_NV_ay=\chi_NV_aF_{\eta,y}\).
It is trace class on \(H_{a,s}\). Also
\(N_a(g)=C_gM_\phi:H\to H_{a,s}\) is trace class:
its normalized physical kernel is precisely (MA-MO.332),
whose compact difference and exponential left decay
were proved there. These norm-coordinate statements do
not depend on the number of places of the fixed field.

The first actual smoothed row of (MA-MO.479) is now
\[
        \mathscr R_{\eta,a}(g)y
          =K_{\eta,a}(g)y-N_a(g)y+O_{\eta,a}(g)Ry.
 \tag{MA-MO.482}
\]
To verify it, use (MA-MO.480) on \(H_{\eta,y}\) and
\(Z_{\eta,a}M_qJ_{\eta,a}-1
   =(M_\phi-Z_{\eta,a}M_\phi J_{\eta,a})-M_\phi\).
Every term is trace class from \(H\) to \(H_{a,s}\).

The second row is
\(R\mathscr R_{\bar\eta,a}(g^\vee)R\), with
\(g^\vee(x)=g(-x)\). The actual Fourier relation in
(MA-MO.479), its scaling by \(g^\vee\), and zero-singular
Poisson summation show that the two rows agree as functions
for every \(y\in\mathcal E\). The paired extension has
values in \(H_{a,s}\oplus H_{-a,s}\), and is trace class
there. Its range is in the closed diagonal, by density of
\(\mathcal E\) in the sum norm. Applying the inverse of
that diagonal isometry gives exactly (MA-MO.473), agreeing
with the original smooth matrix on its whole common sector.
This also proves uniqueness. No unbounded individual
factor has been declared defined on every Hilbert input.

\emph{The actual ideal restriction and quotient action.}
On the dense ideal source \(I_\eta\), the map just
constructed agrees with the original action whose
completion is \(B_{\eta,g}\) in MA-MO-012. Both maps
are bounded, so their restrictions on \(J_\eta\) agree.
In particular \(J_\eta\) is invariant, and its trace is
\(\tau_{\eta,S}(g,\phi)\) by MA-MO-017.

The bounded map \(\mathscr D_{\eta,g}+C_g\) agrees
on \(\mathcal E\) with \(C_gP_+\), whose range is
in \(I_\eta\). Closure therefore puts its entire range
in \(J_\eta\). Since \(C_g\) preserves \(J_\eta\),
the induced quotient map of \(\mathscr D_{\eta,g}\)
is exactly \(-C_g\). In the orthogonal decomposition
\(H=J_\eta\oplus J_\eta^\perp\), the operator is
upper triangular with diagonal blocks \(B_{\eta,g}\)
and the quotient operator \(-C_g\). Compression makes
both blocks trace class. Absolute summation against the
two orthonormal bases gives the trace equality in
(MA-MO.474). Finite-dimensional block compressions converge
in trace norm; their triangular determinant identity and
Fredholm continuity prove the determinant equality as well.

\emph{Full retained zeros and the finite prime evaluation.}
The functional equation between \(\eta\) and \(\bar\eta\),
together with (MA-MO.472) for both, places every completed
zero in \(|\Re\rho-1/2|\le a_*<a\). Thus MA-WP-001--003
retains all multiplicities, with no boundary truncation,
and gives the actual quotient trace
\[
        \operatorname{Tr}_{Q_\eta}C_g
                  =\sum_\rho m_\eta(\rho)G(\rho-1/2).
 \tag{MA-MO.483}
\]
The series is absolutely convergent by the fixed-character
local zero count and rapid Mellin decay. It is the trace
of the actual completed quotient, not an assigned spectral
functional.

For \(g(0)=0\), the full primitive explicit formula in
MA-CH/MT-M30 and Meyer's Theorems~5.8/5.11 reads
\[
       \sum_\rho m_\eta(\rho)G(\rho-1/2)
                   =-\mathcal A_{\infty,\eta}(g)-B_\eta(g).
 \tag{MA-MO.484}
\]
Here \(\mathcal A_{\infty,\eta}\) is the sum of the real
parity and complex gamma integrals, and \(B_\eta\) uses
all primitive finite Euler factors with conjugate coefficients
on the negative argument, as in MA-CH.2 and .6. There is
no principal polar term. The conductor logarithm pairs
with \(g(0)\), hence is zero in this displayed formula.

The actual multiplier \(m_\eta\) contains the local
factors for \(S\). Each unramified finite factor has the
form, with \(q_v=N\mathfrak p_v\),
\[
       \frac{1-\overline{c_{\mathfrak p_v}}q_v^{z-1/2}}
            {1-c_{\mathfrak p_v}q_v^{-z-1/2}}.
 \tag{MA-MO.485}
\]
Its logarithmic derivative on \(\Re z=0\) has the
absolutely convergent expansion
\[
 -\sum_{k\ge1}(\log q_v)q_v^{-k/2}
 \bigl(c_{\mathfrak p_v}^{\,k}e^{-kz\log q_v}
  +\overline{c_{\mathfrak p_v}}^{\,k}e^{kz\log q_v}\bigr).
\]
Mellin inversion therefore contributes the negative of
that place's primitive prime distribution. Ramified finite
factors and different/conductor dilations are nonzero
exponential monomials times phases; their logarithmic
derivatives are constants and pair with \(g(0)\).
The gamma calculation of MA-MO.369, with both signs of
the test argument retained, consequently gives
\[
        \tau_{\eta,S}(g,\phi)
          =-\mathcal A_{\infty,\eta}(g)-B_{\eta,S}(g)
                            \qquad(g(0)=0).
 \tag{MA-MO.486}
\]
Here \(B_{\eta,S}\) sums the finite places in \(S\).
Ramification of \(K/\mathbb Q\) alone does not delete
an unramified character factor. The nonsymmetric residual
also vanishes because it is \(g(0)\int\psi\).
Subtracting (MA-MO.484) from (MA-MO.486) in the actual
trace identity (MA-MO.474) proves (MA-MO.475).
\end{proof}

\paragraph{Conditional scope and actual operators distinguished.}
The half-plane premise is used twice: to identify the
original wide inverse on every smooth common input with
the positive narrow-line inverse, and to retain the entire
completed divisor in (MA-MO.483). The direct forward
theta shift, the local finite phases, and the full explicit
formula themselves require no zero-location assertion.
For a nonprincipal primitive Dirichlet character, A supplies
(MA-MO.472) with \(a_*=0\). For the non-induced cubic
imaginary-quadratic character of MA-CH-003, its own Hecke RH
or the stipulated M supplies that premise; Dirichlet A
alone does not. No claim about its truth is added here.

The full operator \(\mathscr D_{\eta,g}\) is not the
unrestricted assembled block diagonal operator
\(\mathbb T\) of MA-MO-012/019. They agree on the actual
ideal only after the stated identifications. The former has
the negative quotient convolution in (MA-MO.474); the latter
has the separately computed unrestricted trace
\(2\tau_{\eta,S}\). No ideal projector has been dropped
in passing between these statements.

The structural-source locators freshly checked here are
Meyer (2005), supplied \texttt{main.tex} lines~3570--3725
for the sufficiently-large condition and full coinvariant
proof; lines~2630--2675 for the split sequence and
Schwartz-cup inclusion; and \texttt{lem:Hilm\_p\_embed},
\texttt{eq:Poisson\_summation}, and
\texttt{lem:Poisson\_sumpm}, lines~3808--4043, for the
actual matrix. The two-sided local topology and nuclear
return are the independently checked MA-MO-017.
In the inherited cote~1 framework, scans~079--080 and
111--112 through ALPOGE\_BRIDGES BR-TR-01/02, (MA-MO.474)
is an actual nuclear determinant formula along an invariant
Hilbert subspace and quotient. It does not create a positive
full-jet realization, a new motivic weight, or an unconditional
answer to the remaining FT002 arithmetic gate.
